\section{Design}

\subsection{Architectural sketch}

Below is the architectural sketch of the components structured in the diagram that
highlights the software and hardware components.

\begin{figure}[h]
    \centering
    \includegraphics[width=1\textwidth]{./img/architecturediagram.png}
    \caption{Architectural sketch of the project}
    \label{fig:archi}
\end{figure}

This is the description of each of the elements of the diagram:
\begin{enumerate}
  \item \textbf{ESP32 DevKit microcontroller} -- central processing unit running FreeRTOS
    and executing all three application tasks preemptively
  \item \textbf{DHT11 temperature sensor (GPIO\,4)} -- digital sensor communicating over
    a single-wire protocol; read by TaskAcquisition every 2000\,ms
  \item \textbf{Green LED (GPIO\,25)} -- output indicator; ON when temperature alert is OFF
  \item \textbf{Red LED (GPIO\,26)} -- output indicator; ON when temperature alert is active
  \item \textbf{OLED 128$\times$64 SSD1306 (SDA=21, SCL=22, addr 0x3C)} -- primary display
    updated by TaskReporter every 500\,ms with raw, filtered, averaged values and alert state
  \item \textbf{LCD 20$\times$4 I2C (SDA=21, SCL=22, addr 0x27)} -- backup display showing
    averaged temperature and alert state
  \item \textbf{SensorData module} -- shared data store protected by \texttt{xSensorMutex};
    holds raw, filtered, and averaged temperature values, validity flag, alert state,
    sensor selection, and all configurable pipeline parameters
  \item \textbf{TaskAcquisition} -- FreeRTOS task (priority 3) that reads the sensor
    periodically and gives \texttt{xSemNewData}
  \item \textbf{TaskConditioning} -- FreeRTOS task (priority 2) blocked on
    \texttt{xSemNewData}; applies saturation, median filter, and EMA; evaluates alert
    thresholds; updates LEDs
  \item \textbf{TaskReporter} -- FreeRTOS task (priority 1) that updates OLED, LCD, and
    prints STDIO reports at 500\,ms intervals
  \item \textbf{GpioDriver (MCAL)} -- wraps all \texttt{digitalRead}/\texttt{digitalWrite}
    calls; only layer that touches Arduino hardware APIs directly
\end{enumerate}

\subsection{Electrical sketch}

The circuit diagram outlines the setup for this project, showing the connections between
the ESP32, two LEDs (green, red), the DHT11 temperature sensor, the OLED and LCD I2C
modules, and the computer for serial communication. Each LED is connected to its
respective GPIO pin through a $220\,\Omega$ current-limiting resistor. The DHT11 data
line (GPIO\,4) requires a $10\,\mathrm{k\Omega}$ pull-up resistor to 3.3\,V for reliable
communication. Both displays share the I2C bus (SDA=GPIO\,21, SCL=GPIO\,22) at different
addresses (OLED at 0x3C, LCD at 0x27).

\begin{figure}[h]
    \centering
    \includegraphics[width=0.4\textwidth]{./img/circuit_sketch_wokwi.jpg}
    \caption{Circuit sketched on wokwi}
    \label{fig:wokwi}
\end{figure}

Figure \ref{fig:wokwi} details the connections and components that make up the system.
This visual documentation is crucial for replicating the setup or diagnosing issues, as
it clearly shows how each component is integrated into the overall design.

\subsection{Flow chart}

The Figure \ref{fig:flow} illustrates the FreeRTOS task interaction flow. TaskAcquisition
runs every 2000\,ms, reads the DHT11, updates shared data under the mutex, and gives the
binary semaphore. TaskConditioning unblocks immediately, reads the shared raw temperature,
applies the three-stage conditioning pipeline (saturation $\to$ median filter $\to$ EMA),
stores intermediate results, evaluates alert thresholds, and updates the LEDs.
TaskReporter runs every 500\,ms independently and reads shared data to refresh the OLED,
LCD, and serial output.

\begin{figure}[h]
    \centering
    \includegraphics[width=0.8\textwidth]{./img/flow.png}
    \caption{Flow chart of the algorithm}
    \label{fig:flow}
\end{figure}

\clearpage
\subsection{Code design}

Code is developed following the principle of separating the interface from the
implementation, by creating for each header file \texttt{.h} a respective \texttt{.cpp}
file. The project uses \textbf{PlatformIO} with three library directories:

\begin{enumerate}
  \item \textbf{MCAL layer} (\texttt{lib/MCAL/GpioDriver}) -- hardware-specific code that
    wraps \texttt{Arduino.h} GPIO calls. All direct hardware access is encapsulated here.
  \item \textbf{ECAL layer} (\texttt{lib/ECAL/}) -- electronic component drivers:
    \texttt{TempSensor} (DHT11 via Adafruit DHT library),
    \texttt{LedDriver} (two LEDs),
    \texttt{OledDisplay} (Adafruit SSD1306 128$\times$64),
    \texttt{LcdDisplay} (LiquidCrystal\_I2C, backup), and
    \texttt{SerialIO} (UART initialization).
    Each driver header declares its connection pins for easy circuit reference.
  \item \textbf{SRV layer} (\texttt{lib/SRV/}) -- the \texttt{SensorData} shared-state
    module and the three FreeRTOS task implementations.
\end{enumerate}

The \texttt{SensorData} module centralizes all shared data behind mutex-protected
getters and setters, and owns the FreeRTOS synchronization primitives:

\begin{lstlisting}
// Declared in SensorData.h, initialized in SensorData_Init()
extern SemaphoreHandle_t xSemNewData;   // binary semaphore: new data available
extern SemaphoreHandle_t xSensorMutex;  // mutex: protects all shared sensor fields
\end{lstlisting}

The signal conditioning logic is implemented as a three-stage pipeline in
\texttt{TaskConditioning}:
\begin{lstlisting}
/* 1. Saturation -- clamp to physical sensor range */
static float Saturate(float val)
{
    if (val < SATURATION_LOW)  return SATURATION_LOW;
    if (val > SATURATION_HIGH) return SATURATION_HIGH;
    return val;
}

/* 2. Median filter -- sort window, return middle element */
static float MedianFilter_Get(void)
{
    float sorted[MEDIAN_WINDOW];
    uint8_t n = medianFilled;
    for (uint8_t i = 0; i < n; i++) sorted[i] = medianBuf[i];
    // Insertion sort (n <= 5)
    for (uint8_t i = 1; i < n; i++) {
        float key = sorted[i];
        int j = i - 1;
        while (j >= 0 && sorted[j] > key) {
            sorted[j + 1] = sorted[j]; j--;
        }
        sorted[j + 1] = key;
    }
    return sorted[n / 2];
}

/* 3. Exponential Moving Average */
static float EMA_Update(float newVal)
{
    if (!emaInitialised) {
        emaValue = newVal;
        emaInitialised = 1;
    } else {
        emaValue = EMA_ALPHA * newVal
                 + (1.0f - EMA_ALPHA) * emaValue;
    }
    return emaValue;
}
\end{lstlisting}

TaskReporter uses \texttt{vTaskDelayUntil} for drift-corrected 500\,ms periodicity
and outputs data in the \texttt{>name:value} format compatible with the VS Code
Serial Plotter extension, alongside a full structured report every 5\,seconds:
\begin{lstlisting}
// Serial Plotter line (every 500 ms)
Serial.print(">");
Serial.print("raw:"); Serial.print(raw, 2);
Serial.print(",median:"); Serial.print(filtered, 2);
Serial.print(",ema:"); Serial.print(averaged, 2);
Serial.print(",thresh_high:"); Serial.print(TEMP_THRESH_HIGH, 1);
Serial.print(",thresh_low:"); Serial.print(TEMP_THRESH_LOW, 1);
Serial.println();

// Full structured report (every 5 s)
if ((reportCounter++ % FULL_REPORT_EVERY) == 0) {
    printf("\n===== SENSOR REPORT (t=%lu ms) =====\n", ...);
    printf("Raw           : %.2f C\n", raw);
    printf("Saturated+Med : %.2f C  (median window=%d)\n",
           filtered, MEDIAN_WINDOW);
    printf("EMA averaged  : %.2f C  (alpha=%.2f)\n",
           averaged, EMA_ALPHA);
    printf("Alert         : %s  (ON>=%.1f OFF<=%.1f)\n", ...);
    printf("====================================\n\n");
}
\end{lstlisting}
